\documentclass[11pt,a4paper]{article}
\usepackage[T1]{fontenc}
\usepackage[utf8]{inputenc}
\usepackage[french]{babel}
\usepackage{lmodern,microtype}
\usepackage[margin=23mm]{geometry}
\usepackage{booktabs,array,tabularx,longtable,float}
\usepackage{xcolor}
\PassOptionsToPackage{obeyspaces}{url}
\usepackage{hyperref}
\hypersetup{colorlinks=true,linkcolor=black,citecolor=black,urlcolor=blue,pdftitle={Campagne reelle LANCE Qwen UMONS S1--S12 : rapport intermediaire apres fin du lot initial},pdfauthor={Projet LANCE}}
\urlstyle{same}
\setlength{\emergencystretch}{2em}
\setlength{\parskip}{4pt}
\clubpenalty=10000
\widowpenalty=10000
\def\UrlBreaks{\do\/\do\-\do\_\do\.\do\?\do\&\do\=\do\:\do\,}
\newcolumntype{Y}{>{\raggedright\arraybackslash}X}
\newcommand{\code}[1]{{\urlstyle{tt}\nolinkurl{#1}}}
\title{Campagne r\'eelle LANCE Qwen UMONS S1--S12\\\large Rapport scientifique interm\'ediaire apr\`es fin du lot initial}
\author{Projet LANCE\\\small Muse (r\'edaction) ; Codex (relecture et corrections)}
\date{Derni\`ere r\'evision : 5 octobre 2026, 11:48\\Gel des donn\'ees : \code{authoritative-evidence.json} du 5 octobre 2026, 11:32\\Europe/Brussels (UTC+02:00)\\\textbf{Statut : r\'esultats interm\'ediaires ; lot initial termin\'e ; validations apr\`es corrections \`a ex\'ecuter}}
\begin{document}
\maketitle
\begin{abstract}
Ce rapport intermédiaire analyse les premiers runs réels LANCE sur S1--S12, exécutés sur nato/pve-nato avec \code{ollama-umons / qwen3.8:27b}. Le faisceau gelé le 5 octobre 2026 à 11:32 (Europe/Brussels, UTC+02:00) couvre 14 archives et 1\,582 fichiers aux empreintes vérifiées. Sous \code{strict-v3.14 / evidence-v13}, S1 commun et S2--S11 ont un rapport et un score officiels ; S12 échoue en reconnaissance sans score ni rapport, ce qui ne vaut pas zéro. S1 commun obtient 11 VP, 2 FP et 1 FN (F1 confirmé 0,880) ; ses deux diagnostics antérieurs restent séparés. Les revues distinguent cohérence arithmétique, admission des preuves et validité sémantique : des compteurs périmés, des preuves natives rejetées, des surinterprétations et des fixtures défectueuses motivent quatorze corrections. Leurs tests logiciels et rejeux hors réseau ne démontrent aucun gain expérimental. Au gel, leur CI finale et leur déploiement sont en attente ; leur validation réelle reste à faire. Les revues ciblées sont achevées, mais les faux négatifs et preuves secondaires ne sont pas exhaustivement certifiés. L'arrêt demandé après S12 a reçu HTTP 401 ; le lot a terminé naturellement, S13--S19 inclus, dont les scores sont exclus. Le bilan reste une campagne de développement, sans évaluation indépendante ni conclusion globale de fiabilité.
\end{abstract}
\clearpage
\setcounter{tocdepth}{2}
{\small\setlength{\parskip}{0pt}\tableofcontents}
\clearpage
\section{Statut du document et r\`egles de lecture}
\label{sec:statut}
Ce document est un \textbf{rapport scientifique interm\'ediaire}, mis \`a jour apr\`es la fin du lot initial. Il \'etend le snapshot pr\'ec\'edent (gel 2026-10-05 10:12, r\'evision 10:36) sans r\'e\'ecrire ses r\'esultats historiques. La date de derni\`ere r\'evision identifie cette version de r\'edaction ; elle ne constitue pas une validation des corrections ni une certification des preuves.

R\`egles appliqu\'ees dans tout le rapport :
\begin{itemize}
\item Les chiffres des scores archiv\'es sont repris tels quels, sans r\'e\'evaluation. Les co\^uts affich\'es sont des \textbf{estimations applicatives} (\code{cost_is_estimate: true}), ni facture ni tarification UMONS.
\item Le \textbf{F1 final} est toujours celui de la population \textbf{confirm\'ee} (\code{official_funnel.confirmed}). Les populations candidates et filtr\'ees sont des instantan\'es distincts et ne s'y substituent pas.
\item \textbf{S12 n'a pas de score} : sa reconnaissance a \'echou\'e et l'endpoint officiel r\'epond 404. Il n'est jamais rendu comme un z\'ero.
\item Les deux S1 de diagnostic forment un groupe \textbf{diagnostic} s\'epar\'e du S1 commun et du lot : aucune moyenne ni agr\'egation ne les fusionne.
\item Au gel, les quatorze corrections préparées sont \textbf{pouss\'ees, CI en succ\`es ou en cours selon le commit, non d\'eploy\'ees} : leurs rejeux hors r\'eseau et leurs tests logiciels ne sont pas des validations r\'eelles et ne modifient ni la v\'erit\'e terrain ni les scores historiques.
\item Les revues S8--S11 sont \textbf{cibl\'ees et non exhaustives} : exclusions natives et grandes surinterpr\'etations trait\'ees, revue compl\`ete des faux n\'egatifs, preuves secondaires et validit\'e de la v\'erit\'e terrain encore ouverte. Cette limite est signal\'ee et non compens\'ee par hypoth\`ese.
\item Les faits gel\'es du snapshot priment sur toute donn\'ee live observ\'ee apr\`es le gel.
\end{itemize}
Rédaction : Muse, en lecture seule sur archives ; relecture des chiffres, correction et application de la source : Codex. Le rendu de cette nouvelle version est contrôlé séparément des relectures antérieures ; son périmètre est consigné dans le README du sujet.

\section{Objectif et p\'erim\`etre}
\label{sec:objectif}
L'objectif est de fiabiliser la cha\^ine LANCE existante avant son utilisation par une doctorante et des \'etudiants (travaux de fin d'\'etudes) : pr\'eparation du laboratoire, six phases d'audit, preuves trac\'ees, m\'etriques coh\'erentes, affichage des r\'esultats et nettoyage. Les corrections sont motiv\'ees par les sorties r\'eelles et compl\'et\'ees par des v\'erifications cibl\'ees du code.

Le p\'erim\`etre demand\'e est \textbf{S1--S12, six phases, profil \code{full}}, sans jeu tenu \`a l'\'ecart de la mise au point : aucune \'evaluation ind\'ependante sur donn\'ees s\'equestr\'ees n'est revendiqu\'ee. Il n'y a ni nouveau lot S13--S29 ni jeu \code{held-out}. Le plan d'étude initial allait jusqu'à S29 (S1--S19 en développement, S20--S29 en évaluation), mais le seul lot initial effectivement lancé contient S2--S19 puis S1. Le périmètre de l'étude a été réduit à S1--S12 pendant son exécution ; la requ\^ete d'origine reste archiv\'ee pour tra\c{c}abilit\'e. Les variantes historiques S1h et S4h restent exclues. Aucun audit local du benchmark ni aucune intervention sur le homelab personnel ne rel\`event de ce rapport.

\'Etat au gel : lot initial \textbf{termin\'e} avec nettoyage ; deux S1 de diagnostic archiv\'es et s\'epar\'es ; S1 commun du lot termin\'e ; S2--S11 avec audits, rapports et scores officiels archiv\'es ; S12 en \'echec de reconnaissance sans score ni rapport ; validations r\'eelles apr\`es corrections \`a ex\'ecuter. Les d\'emarrages hors p\'erim\`etre S13--S19 rel\`event de l'incident d'arr\^et d\'ecrit en section~\ref{sec:scope} et leurs scores sont exclus des r\'esultats.

\section{M\'ethode et conditions de reproduction}
\label{sec:methode}
\subsection{Cible, fournisseur, profil et versions}
Les sc\'enarios sont d\'eploy\'es sur \code{pve-nato} par la plateforme h\'eberg\'ee sur \code{nato-master}. Le homelab personnel et l'ex\'ecution locale des sc\'enarios sont exclus. Les demandes passent par l'API \url{http://100.103.253.86:8501} ; les journaux sont ensuite archivés pour analyse. Muse a rédigé en lecture seule sur le faisceau gelé. Codex a séparément inspecté les preuves et validé les corrections logicielles ; les audits réels restent exécutés sur NATO.

Le fournisseur est \code{ollama-umons} (Universit\'e de Mons) et le mod\`ele est \code{qwen3.8:27b}. Le profil demand\'e est \code{full} : les six phases sont activ\'ees (analyse du graphe, reconnaissance, analyse des vuln\'erabilit\'es par machine, v\'erification par preuves, \'evaluation des acc\`es, rapport) avec nettoyage automatique. La topologie publique est accessible au pipeline ; la v\'erit\'e terrain reste r\'eserv\'ee \`a l'\'evaluateur. Le tag du mod\`ele ne constitue pas une version immuable et son empreinte binaire n'est pas consign\'ee dans chaque run. Les conditions de r\'ef\'erence et d'ex\'ecution sont d\'ecrites dans les sources~\cite{catalogue,contrat}.

Identifiants abr\'eg\'es des commits archiv\'es avec chaque run : \code{7ac478c} (premier S1 diagnostic), \code{457d060} (second S1 diagnostic), \code{e9d7d5590d4adf7a3675ae49d95fe8926127f790} (S2--S12 et S1 commun du lot \code{batch-aeedb60283ef}). Les deux diagnostics S1 portent un plafond applicatif de co\^ut estim\'e de 2~dollars par run, absent du lot. Le d\'elai de phase~3 par appareil est de 240~secondes ; ces \'ech\'eances ne sont pas attribu\'ees \`a une panne r\'eseau. Les nettoyages des deux diagnostics S1, des runs S2--S12 et du S1 commun sont \code{completed}. S12 est un audit \code{failed} ; son nettoyage r\'eussi ne transforme pas cet \'echec en audit termin\'e.

Les requêtes \code{restored-s1-request.json}, \code{corrected-s1-request.json} et \code{development-batch-request.json} sont conservées sous le dossier d'artefacts cité ci-dessous. Les configurations effectives figurent dans \code{run_meta.json}, champ \code{execution_profile_config} : plafonds de 20/50/10/10/80 tours pour les phases 1--5, un tour de rapport, respectivement 4\,096 jetons par réponse pour les phases 1--4, 16\,384 pour l'intrusion, 2\,048 pour le rapport. Les boucles locales ont des limites séparées consignées dans ce même champ ; le lot n'a pas de plafond monétaire. La plateforme API conserve l'\'etat d'ex\'ecution ; le lot ne fige ni les poids du mod\`ele ni un service UMONS immuable. Un rerun doit conserver ces diff\'erences de configuration dans ses m\'etadonn\'ees.

\subsection{Entonnoir de confirmation et contrats}
Dans l'entonnoir final confirm\'e, \textbf{VP} (vrais positifs) désigne une déclaration confirmée, appariée à une référence attendue et soutenue par une preuve admise ; une référence ne reçoit qu'un crédit, \textbf{FP} (faux positifs) une d\'eclaration confirm\'ee non cr\'edit\'ee par la r\'ef\'erence, \textbf{FN} (faux n\'egatifs) une référence attendue sans crédit de confirmation admissible. Avec \(P\) pour la pr\'ecision et \(R\) pour le rappel, \(P=VP/(VP+FP)\), \(R=VP/(VP+FN)\) et \(F_1=2PR/(P+R)\), calcul\'ees sur l'entonnoir confirm\'e. Des indicateurs historiques pond\'er\'es pr\'esents dans \code{score.json} utilisent d'autres populations et ne sont ni interchangeables avec ce F1 ni substitu\'es ici.

Les \'etapes \textbf{candidates}, \textbf{filtr\'ees} et \textbf{confirm\'ees} sont trois instantan\'es \'evalu\'es s\'epar\'ement : registre brut de phase~3, projection canonique apr\`es filtrage, d\'eclarations confirm\'ees de phase~4. Chaque population est d\'edupliqu\'ee suivant le contrat structurel ; ce regroupement conservateur peut laisser des d\'eclarations distinctes visant une m\^eme r\'ef\'erence, que l'appariement ne cr\'edite qu'une fois (autres class\'ees redondantes). Les FP incluent des confirmations soutenues hors r\'ef\'erence et ces redondances ; ils ne sont par d\'efaut ni qualifi\'es d'hallucinations ni compt\'es comme vrais positifs ajout\'es. Les cat\'egories de revendications (\code{insufficient_evidence}, \code{supported_outside_reference}, \code{redundant_reference_claim}, \code{reference_assignment_conflict}, \code{contradicted}) sont rapport\'ees telles quelles : une exclusion g\'en\'erique pour preuve insuffisante n'est pas automatiquement un faux n\'egatif.

Tous les runs du gel utilisent \code{strict-v3.14} / \code{evidence-v13}. Le contrat \code{strict-v3.15} / \code{evidence-v14} est une \'etape pr\'epar\'ee interm\'ediaire \textbf{jamais d\'eploy\'ee} ; le contrat pr\'epar\'e final pour les futurs runs correctifs est \code{strict-v3.16} / \code{evidence-v15} (preuve propre par ressource, lectures natives typ\'ees, gardes robots/HTTP). Aucun effet de ces contrats pr\'epar\'es ne peut \^etre pr\'esent\'e comme un gain \`a contrat constant, et aucune validation r\'eelle corrig\'ee n'existe au gel.

\subsection{Contr\^oles r\'ealis\'es et artefacts}
Pour chaque run sont rapproch\'es : journaux de pr\'eparation, livrables de phases, registre des outils, traces fournisseur, r\'esum\'e de consommation, m\'etriques officielles (\code{score.json}, \code{run_meta.json}, \code{cost_summary.json}, \code{metric-audit.json}), revues et rejeux du dossier parent. Les erreurs de transport sont distingu\'ees des r\'eponses n\'egatives des services et des preuves insuffisantes. Une absence de r\'esultat reste une absence. Z\'ero saut v\'erifi\'e ne prouve pas l'absence de possibilit\'e : les tentatives de pivot sont indisponibles, pas nulles. Un pr\'e-contr\^ole \code{preflight-v1} r\'eussi n'atteste pas toutes les propri\'et\'es de la v\'erit\'e terrain. Un contr\^ole arithm\'etique et des r\'ef\'erences ne suffisent pas \`a prouver la validit\'e s\'emantique. Le code, les tests logiciels et les rejeux hors r\'eseau ne prouvent pas un gain exp\'erimental.

Les artefacts primaires sont conservés sous \path{output/validation/2026-10-04-qwen-umons/} : requ\^etes, journaux, archives par run, scores officiels, manifestes et empreintes \code{collection-manifest.json}. Le faisceau de données de cette version est \path{muse-baseline-report-2026-10-05/authoritative-evidence.json}, sous ce dossier (r\'evision 2026-10-05 11:32 Europe/Brussels). Les quatorze archives couvertes repr\'esentent 1\,582 fichiers dont les empreintes ont \'et\'e contr\^ol\'ees, sans r\'e\'ecriture des m\'etriques originales.

Les revues et rejeux sont cit\'es aupr\`es des constats concern\'es et regroup\'es dans les r\'ef\'erences~\cite{artefacts,revues}. L'incident d'arr\^et est consign\'e avec sa politique et son \'etat d'observation ; la validation logicielle des corrections figure dans les artefacts CI d\'edi\'es. Les journaux originaux et leurs empreintes restent les pi\`eces de r\'ef\'erence.

\section{R\'esultats S1--S12 au gel}
\label{sec:resultats}
\subsection{\'Etat d'avancement}
Le tableau~\ref{tab:etat} r\'esume le statut de chaque sc\'enario. Il se lit comme un \'etat d'avancement, pas comme un palmar\`es : les r\'esultats d\'ecrivent des sc\'enarios, r\'ef\'erences, versions de code et configurations distincts, et quelques r\'ealisations S1 ne suffisent pas \`a \'etablir un gain statistique.
\begin{table}[H]
\centering
\caption{\'Etat S1--S12 au gel (2026-10-05 11:32 Europe/Brussels). Lot initial termin\'e.}
\label{tab:etat}
\begin{tabularx}{\linewidth}{l Y}
\toprule
Sc\'enario & \'Etat au gel\\
\midrule
S1 diagnostiques & Deux runs ant\'erieurs archiv\'es (partiel puis termin\'e) ; groupe diagnostic, non agr\'eg\'es au lot ni au S1 commun.\\
S1 commun & Run \code{103656} termin\'e : six phases \code{completed}, nettoyage r\'eussi, F1 confirm\'e 0,880.\\
S2--S11 & Audits et rapports pr\'esents, nettoyages r\'eussis, scores officiels archiv\'es ; statuts \code{partial} sauf S6 \code{completed} (voir tableaux~\ref{tab:funnel} et~\ref{tab:conso}).\\
S12 & Reconnaissance incompl\`ete (\code{failed}), phases 3--6 \code{skipped}, score indisponible (HTTP 404), pas de rapport final. N'est pas un z\'ero.\\
\bottomrule
\end{tabularx}
\end{table}
Le lot \code{batch-aeedb60283ef} a d\'emarr\'e le 4 octobre 2026 sur \code{e9d7d55}. Sa requ\^ete d'origine est conserv\'ee pour tra\c{c}abilit\'e. L'arr\^et demand\'e apr\`es S12 a \'echou\'e sur l'authentification (voir section~\ref{sec:scope}) ; le lot a ensuite roul\'e jusqu'\`a son terme naturel, S1 commun inclus, avec nettoyage.
\subsection{Entonnoir final confirm\'e officiel}
Le tableau~\ref{tab:funnel} donne l'entonnoir confirm\'e officiel de chaque run archiv\'e. Les fractions P/R/F1 utilisent les VP/FP/FN affich\'es. Tous les runs sont en \code{strict-v3.14} / \code{evidence-v13}. S12 sans score est absent du tableau chiffr\'e par construction.
\begin{table}[H]
\centering
\small
\caption{Entonnoir final confirm\'e officiel. S12 sans score : absent par construction.}
\label{tab:funnel}
\setlength{\tabcolsep}{4pt}
\begin{tabular}{llllrrr}
\toprule
Run & Sc\'en. / groupe & Commit & Statut & VP/FP/FN & P / R & F1\\
\midrule
\code{153134} & S1 diagn. & \code{7ac478c} & partiel & 11/1/1 & 0,917 / 0,917 & 0,917\\
\code{162046} & S1 diagn. & \code{457d060} & termin\'e & 11/3/1 & 0,786 / 0,917 & 0,846\\
\code{170828} & S2 d\'ev. & \code{e9d7d55} & partiel & 12/7/1 & 0,632 / 0,923 & 0,750\\
\code{180449} & S3 d\'ev. & \code{e9d7d55} & partiel & 12/14/6 & 0,462 / 0,667 & 0,545\\
\code{191627} & S4 d\'ev. & \code{e9d7d55} & partiel & 9/13/9 & 0,409 / 0,500 & 0,450\\
\code{202341} & S5 d\'ev. & \code{e9d7d55} & partiel & 10/12/5 & 0,455 / 0,667 & 0,541\\
\code{212837} & S6 d\'ev. & \code{e9d7d55} & termin\'e & 12/7/4 & 0,632 / 0,750 & 0,686\\
\code{222716} & S7 d\'ev. & \code{e9d7d55} & partiel & 8/7/6 & 0,533 / 0,571 & 0,552\\
\code{232009} & S8 d\'ev. & \code{e9d7d55} & partiel & 7/8/7 & 0,467 / 0,500 & 0,483\\
\code{003954} & S9 d\'ev. & \code{e9d7d55} & partiel & 7/8/4 & 0,467 / 0,636 & 0,538\\
\code{015017} & S10 d\'ev. & \code{e9d7d55} & partiel & 3/2/10 & 0,600 / 0,231 & 0,333\\
\code{023056} & S11 d\'ev. & \code{e9d7d55} & partiel & 13/18/10 & 0,419 / 0,565 & 0,481\\
\code{103656} & S1 commun & \code{e9d7d55} & termin\'e & 11/2/1 & 0,846 / 0,917 & 0,880\\
\bottomrule
\end{tabular}
\end{table}
Les identifiants de run sont abr\'eg\'es : pr\'efixe \code{2026-10-04} pour les deux S1 diagnostics et S2--S8, puis \code{2026-10-05} pour S9--S12 et le S1 commun. Les dossiers complets permettent de retrouver les archives.

R\'ef\'erences attendues et appariements : S1 diagnostics 12 r\'ef. (11 appari\'ees chacun) ; S2 13 (12) ; S3 18 (12) ; S4 18 (9) ; S5 15 (10) ; S6 16 (12) ; S7 14 (8) ; S8 14 (9 appariements, 7 VP) ; S9 11 (8 appariements, 7 VP) ; S10 13 (4 appariements, 3 VP) ; S11 23 (16 appariements, 13 VP) ; S1 commun 12 (11 appari\'ees). Le nombre d'appariements est un diagnostic structurel distinct du cr\'edit VP : une d\'eclaration peut correspondre \`a une r\'ef\'erence attendue sans disposer d'une preuve admise par le contrat final. Les revendications hors r\'ef\'erence et redondantes suivent la nomenclature du snapshot : par exemple S2 6 soutenues hors r\'ef\'erence et 1 redondante ; S3 12 et 2 ; S4 10 et 3 ; S5 10 et 2 ; S6 5 et 2 ; S7 5 et 2 ; S8 4 et 2 (plus 2 preuves insuffisantes) ; S9 4 et 2 (plus 2 insuffisantes) ; S10 1 et 0 (plus 1 insuffisante) ; S11 11 et 3 (plus 4 insuffisantes) ; S1 commun 1 soutenue hors r\'ef\'erence et 1 redondante. Aucune revendication contredite ni conflit d'assignation n'est enregistr\'e sur S1--S11 et le S1 commun.
\subsection{Consommation, acc\`es et qualit\'e d'ex\'ecution}
Le tableau~\ref{tab:conso} rapporte dur\'ees r\'eelles, tours mod\`ele, appels et erreurs d'outils, acc\`es corrobor\'es sur cibles attendues, co\^uts estim\'es et contr\^oles arithm\'etiques. Les accès comptent des cibles avec accès de service corroboré, y compris anonyme ; ils ne comptent pas des shells ni des pivots : \textbf{z\'ero saut causal v\'erifi\'e} sur tous les runs (tentatives de pivot indisponibles).
\begin{table}[H]
\centering
\small
\caption{Consommation et qualit\'e d'ex\'ecution. Co\^uts applicatifs estim\'es en USD. S12 : reconnaissance seule.}
\label{tab:conso}
\setlength{\tabcolsep}{4pt}
\begin{tabular}{lrrrrrc}
\toprule
Run & Dur\'ee (s) & Tours & \shortstack{Appels /\\erreurs} & \shortstack{Acc\`es /\\cibles} & \shortstack{Co\^ut\\estim\'e} & \shortstack{Contr\^oles\\35 checks}\\
\midrule
\code{153134} (S1) & 2\,404,6 & 111 & 211 / 2 & 2 / 4 & 1,7801 & 35/35\\
\code{162046} (S1) & 2\,254,6 & 93 & 205 / 5 & 2 / 4 & 1,3979 & 35/35\\
\code{170828} (S2) & 3\,353,0 & 149 & 304 / 8 & 2 / 6 & 2,2667 & 33/35\\
\code{180449} (S3) & 4\,258,5 & 210 & 368 / 10 & 3 / 8 & 3,2621 & 32/35\\
\code{191627} (S4) & 3\,995,5 & 178 & 354 / 10 & 2 / 8 & 2,8064 & 35/35\\
\code{202341} (S5) & 3\,858,5 & 170 & 389 / 5 & 2 / 8 & 2,6637 & 35/35\\
\code{212837} (S6) & 3\,490,1 & 169 & 357 / 2 & 1 / 6 & 2,5875 & 32/35\\
\code{222716} (S7) & 3\,143,9 & 139 & 303 / 8 & 2 / 6 & 2,1671 & 35/35\\
\code{232009} (S8) & 4\,746,5 & 210 & 385 / 36 & 2 / 8 & 3,3760 & 33/35\\
\code{003954} (S9) & 4\,194,6 & 161 & 239 / 30 & 2 / 6 & 2,6876 & 35/35\\
\code{015017} (S10) & 2\,410,9 & 92 & 259 / 49 & 2 / 6 & 1,4379 & 35/35\\
\code{023056} (S11) & 7\,351,0 & 297 & 555 / 53 & 3 / 15 & 4,7012 & 35/35\\
\code{043439} (S12) & 1\,371,5 & 11 & 23 / 0 & --- & 0,2174 & sans objet\\
\code{103656} (S1) & 2\,161,6 & 102 & 207 / 5 & 2 / 4 & 1,6661 & 33/35\\
\bottomrule
\end{tabular}
\end{table}
V\'erifications archiv\'ees (lignes) : S1 diagnostics 18 et 17 ; S2 30 ; S3 43 ; S4 32 ; S5 34 ; S6 36 ; S7 24 ; S8 31 ; S9 31 ; S10 14 ; S11 54 ; S12 0 ; S1 commun 17 lignes (16 testées, une ignorée dans la projection, 13 confirmées, 3 inconclusives). Empreintes d'archives v\'erifi\'ees : de 82 \`a 204 hachages selon les runs (S12 : 16), total 1\,582 fichiers. Les \'ecarts arithm\'etiques sont des d\'efauts de projection ou de compteurs, pas des scores alternatifs : S2 (candidats 29 observ\'es contre 30, ignor\'es 0 contre 1), S3 (dont \code{VULN-013} perdue de la projection), S6 (\code{VULN-014} et \code{VULN-016} perdues), S8 (30 contre 31, 1 ignor\'e manquant), S1 commun (compteur de candidats 16 contre 17 lignes, ignorés 0 dans le résumé contre 1 dans la projection). S4, S5, S7, S9--S11 passent 35/35, ce qui atteste le calcul, pas la s\'emantique.
\subsection{S1 commun du lot et S1 diagnostiques}
Le run \code{2026-10-05_103656} (\code{e9d7d55}) est le \textbf{S1 commun} du lot : six phases \code{completed}, nettoyage \code{completed}, 13 d\'eclarations confirm\'ees pour 11 VP, 2 FP, 1 FN, F1 0,880. La projection contient 17 lignes : 16 vérifications effectuées et une \code{SKIPPED}, avec 13 confirmations, 3 inconclusives et 0 erreur. Le résumé conserve deux compteurs erronés : \code{candidate_count}=16 au lieu de 17 et \code{skipped_count}=0 au lieu de 1 ; le nombre de tests effectués (16) est correct. Co\^ut estim\'e 1,6661~USD pour 207 appels, 5 erreurs et 102 tours ; 2 acc\`es corrobor\'es sur 4 cibles, z\'ero saut v\'erifi\'e. La revue d\'edi\'ee \code{s1-common-baseline-review.json} a relu 13 confirmations, 13 descriptions et 22 r\'ef\'erences de preuves : l'avertissement scanner DHEat ne d\'emontre pas un d\'eni de service ; la configuration SSH offrant un chiffrement faible atteste une configuration offerte, pas une attaque de l'homme du milieu ; \code{admin/admin} (uid 1000) prouve une session du compte \code{admin}, sans preuve de ses privilèges administrateurs, de root ni d'un défaut constructeur ; la base \code{.14} annonc\'ee dans la configuration applicative n'est ni dans la topologie ni test\'ee (pas de pivot) ; le canal WS poss\`ede des \code{CONNACK}, \code{SUBACK} et \code{PUBLISH} r\'eels, au-del\`a du premier constat mais avec messages limit\'es ; l'acc\`es routeur root avec \code{shadow} vide et les succ\`es \code{root/root} et \code{root/toor} prouvent l'acc\`es root sans prouver chaque mot de passe. La revue s\'emantique d\'etaill\'ee reste partielle : pas de certification exhaustive des faux n\'egatifs, preuves secondaires ou fixtures.

Les deux runs diagnostiques ant\'erieurs restent \textbf{toujours s\'epar\'es} : \code{153134} (\code{7ac478c}, partiel, phase~3 avec erreurs de workers, F1 0,917) et \code{162046} (\code{457d060}, termin\'e avec graphe r\'ecup\'er\'e, F1 0,846 avec deux confirmations redondantes suppl\'ementaires et la m\^eme confirmation SSH hors r\'ef\'erence). La baisse du F1 entre les deux diagnostics ne correspond pas \`a une perte de rappel (11 VP et 1 FN dans les deux cas). La r\'ef\'erence manqu\'ee commune est le bandeau SSH. Aucune moyenne ni agr\'egation n'est calcul\'ee entre ces trois r\'ealisations S1.

\subsection{S2 et S3}
S2 (\code{170828}) : 19 d\'eclarations confirm\'ees, 10 inconclusives en phase~4, 30 constats canoniques pour 29 effectivement test\'es ; une sonde HTTP sur le port WebSocket 9001 du broker a re\c{c}u un refus de connexion (refus de transport av\'er\'e, pas un faux timeout). Six confirmations soutenues hors r\'ef\'erence et une exposition HTTP redondante ; r\'ef\'erence manqu\'ee : interface d'administration web du routeur c\^ot\'e WAN. Deux acc\`es corrobor\'es sur six cibles, sans pivot causal. La revue note aussi une prose SSH excessive (courbes NIST pr\'esent\'ees comme suspectes) correctement laiss\'ee inconclusive par la politique de preuve.

S3 (\code{180449}) : huit machines analys\'ees ; les deux sondes WebSocket refus\'ees suivaient pourtant une sonde nmap fra\^iche et r\'eussie signalant explicitement \code{9001/tcp closed} (cause commune avec S2). La projection historique ne conserve que 38 test\'ees et 26 confirm\'ees ; \code{VULN-013} (constat MQTT v\'erifi\'e) a disparu de la projection canonique apr\`es une red\'ecouverte redondante d'h\^otes connus. Deux listes de chemins HTTP ont \'et\'e rejet\'ees comme non primaires \`a cause d'une normalisation d\'efectueuse. Trois acc\`es corrobor\'es sur huit cibles, sans pivot v\'erifi\'e.

\subsection{S4 et S5}
S4 (\code{191627}) : huit machines analys\'ees ; 32 v\'erifications (22 confirmations, 10 inconclusives, z\'ero erreur). L'observation Modbus de phase~4 est limit\'ee \`a une identification sans authentification malgr\'e une formulation initiale excessive sur l'\'ecriture de registres ; la phase~5 enregistre s\'epar\'ement une lecture du registre de simulation puis des acquittements d'\'ecriture et de restauration, sans relecture ind\'ependante apr\`es restauration : action protocolaire pr\'esent\'ee s\'epar\'ement, ni shell ni pivot. Pour \code{robots.txt}, le corps frais de phase~4 liste explicitement des chemins candidats \code{Disallow}, mais lister un chemin ne prouve ni son existence ni son contenu sensible ; le v\'erificateur g\'en\'erique limit\'e aux versions laisse cette divulgation non confirm\'ee sous le contrat d\'eploy\'e. Diagnostics erron\'es : alias Nmap \code{mbap} non reconnu (faux diagnostic malgr\'e la sonde Modbus d\'eclar\'ee), adresse IP cliente MySQL du runner prise pour une nouvelle cible, inventaire binaire local refus\'e \`a cause de noms de programmes r\'eseau dans les arguments de recherche.

S5 (\code{202341}) : nettoyage r\'eussi ; six analyses sur huit valid\'ees, routeur et HVAC ayant expir\'e \`a 240 secondes (d\'elai par appareil, pas de panne UMONS \'etablie). 34 v\'erifications (22 confirmations, 12 inconclusives). La seconde cam\'era illustre une limite du contrat : l'acc\`es anonyme confirm\'e ne s'apparie pas \`a la r\'ef\'erence qui n'accepte qu'un autre chemin (non-appariement conforme au contrat conserv\'e) ; deux hypoth\`eses extrapolent un flux r\'eel alors que le script d'injection \'ecrit des octets al\'eatoires et que les corps sont tronqu\'es. Deux acc\`es SSH corrobor\'es sur huit cibles (routeur et NVR), sans pivot ; une commande compos\'ee \'echoue sur \code{whoami} absent du routeur, l'acc\`es reposant sur une autre sonde r\'eussie.

\subsection{S6 et S7}
S6 (\code{212837}) : seul run du lot avec statut \code{completed} outre le S1 commun, six analyses sur six, sans erreur de scanner. La revue documente trois \'ecarts : deux v\'erifications MQTT perdues (\code{VULN-014}, \code{VULN-016}) lors d'une red\'ecouverte d'h\^otes connus (preuves fra\^iches conserv\'ees au registre) et un cr\'edit s\'emantiquement incorrect (une banni\`ere \code{nginx} associ\'ee aux chemins publi\'es dans \code{robots.txt}, sans preuve du contenu de ce fichier). Le score historique (12 cr\'edit\'ees sur 16, F1 0,686) est conserv\'e ; aucun gain combin\'e hypoth\'etique n'est calcul\'e. Phase~5 : un acc\`es routeur corrobor\'e (SSH et Telnet comptant pour une seule cible), z\'ero saut v\'erifi\'e pour 3 attendus.

S7 (\code{222716}) : rapport final pr\'esent malgr\'e le statut partiel, 24 v\'erifications test\'ees et planifi\'ees sans identifiant perdu, 35/35 contr\^oles arithm\'etiques, deux acc\`es sur six cibles, z\'ero saut pour 3 attendus. La revue reproduit les 24 verdicts et r\'ef\'erences avec le m\^eme interpr\'eteur ; cette reproductibilit\'e logicielle ne constitue pas une confirmation s\'emantique ind\'ependante de chaque affirmation. Trois points restent ouverts et motivent des corrections non d\'eploy\'ees : composantes du chemin local \code{.ssh/} faussement reconnues comme commandes r\'eseau, identifiant MySQL \code{root} vide non confirm\'e, et fixture privil\'egi\'ee d\'efectueuse (l'utilitaire en mode 4755 est un script Bash, et Linux ignore les bits set-UID sur les scripts~\cite{execve} ; la lecture prot\'eg\'ee reste refus\'ee). Les refus de l'outil m\'emoire sont intentionnels (m\'emoire \'episodique d\'esactiv\'ee).

\subsection{S8--S11 : scores archiv\'es, revues cibl\'ees termin\'ees, revue exhaustive ouverte}
S8 (\code{232009}), S9 (\code{003954}), S10 (\code{015017}) et S11 (\code{023056}) disposent chacun d'un rapport final archiv\'e, de nettoyages r\'eussis et de scores officiels archiv\'es repris au tableau~\ref{tab:funnel}. Avoir un rapport final ne signifie pas un pipeline irr\'eprochable : S8--S11 restent en statut \code{partial} (phase~3 avec erreurs de workers), et S10 illustre un rappel faible (0,231, 10 FN sur 13 r\'ef\'erences) malgr\'e une pr\'ecision de 0,600.

Niveau de preuve disponible au gel : les revues s\'emantiques cibl\'ees S8, S9, S10 et S11 sont termin\'ees sur les exclusions natives et les grandes surinterpr\'etations, avec rejeu pur d'interpr\'etation (228 confirmations sur 13 archives hash-v\'erifi\'ees rejou\'ees ; exactement 9 modifications d'admission de preuves sur S8--S11, en comparant \code{b39b0b1} préparé/non déployé à \code{0f4e5ff}, et non au baseline \code{e9d7d55} ; vérité terrain non chargée, scores non recalculés). Aucune revue exhaustive des faux n\'egatifs, preuves secondaires ou validit\'e de la v\'erit\'e terrain n'est archivée : les exclusions g\'en\'eriques pour preuve insuffisante (S8 : 2, S9 : 2, S10 : 1, S11 : 4, ce dernier restant FP par contrat) ne sont pas compt\'ees par d\'efaut comme des faux n\'egatifs suppl\'ementaires. Les F1 S8--S11 sont donc rapport\'es comme scores archiv\'es non r\'e\'evalu\'es, pas comme constats s\'emantiquement certifi\'es.

Constats ciblés par sc\'enario (port\'ee restreinte, scores inchang\'es). S8 : lectures Redis (\code{GET} de secrets et \code{PING} sans authentification) r\'eellement archiv\'ees mais rejet\'ees par la liaison historique au faux endpoint HTTP racine ; les références des traces originales sont conservées dans le rejeu pur ; seule leur admission change. S9 : pour \code{VULN-029}, la trace \code{tc-1da51da6ee8a4d9290255f543cee8a16} contient une réponse GET SNMP \code{public} corrélée et lisible ; sa sonde \code{private} distincte expire. Un autre outil, Nmap \code{snmp-brute}, rapporte aussi \code{private} valide. Ces observations différentes ne démontrent aucune écriture SNMP et ne doivent pas être fusionnées en accès lecture/écriture. S10 : lectures FTP au-del\`a des listings (dumps et configuration avec secrets) r\'eellement retourn\'ees par les traces, rejet\'ees par le faux endpoint composite historique. S11 (le plus co\^uteux : 7\,351~s, 297 tours, 555 appels, 53 erreurs, 54 lignes de v\'erification, 204 hachages) : 31 confirmations relues dont 27 preuves admises historiquement ; 13 VP, 18 FP, 10 FN, F1 0,481 ; quatre \code{insufficient_evidence} restent FP. Lectures FTP et Redis r\'eellement archiv\'ees mais exclues par la liaison historique (ne pas inventer de FN ni de gain). OTA : requ\^etes non sign\'ees r\'epondant \code{nginx} statique 200, aucun firmware install\'e ni ex\'ecut\'e. Modbus : identification via Nmap \code{modbus-discover}, pas de lecture/écriture de registres démontrée. HMI : adresse de base annoncée mais incohérente avec la topologie, ni identifiants valid\'es ni pivot ; jeton Redis litt\'eral non utilisable. MySQL : \code{root} vide en TCP et table des utilisateurs lisible, sans preuve d'\'ecriture, suppression ou acc\`es toute adresse. Phase~5 : trois acc\`es, z\'ero saut, vivier brut de 10 dont 5 noms de cl\'es d\'ecoup\'es (ne pas pr\'esenter dix secrets) ; invite Telnet \code{admin/admin} n'attestant pas le mot de passe.

Inspections compl\'ementaires ant\'erieures conserv\'ees : sur S9, le mémo de \code{VULN-029} restreint sa propre conclusion au GET \code{public}, tandis que la projection historique conserve la prose candidate \code{private/écriture} ; le correctif préparé restreint aussi la description canonique ; sur S10, le port Node-RED 1880 refuse les connexions et le script d'injection contient un d\'efaut de syntaxe reproduit hors r\'eseau, correction d\'edi\'ee \`a valider r\'eellement. Ces inspections ne remplacent pas la revue exhaustive restante.

\subsection{S12 : reconnaissance incompl\`ete, pas de z\'ero}
S12 (\code{2026-10-05_043439}) : le graphe est récupéré (\code{completed:recovered}). La reconnaissance échoue (\code{failed}, couverture incomplète ou invalide) ; les phases 3--6 sont omises (\code{skipped:prerequisites}). Co\^ut estim\'e 0,2174~USD pour 23 appels et 11 tours, sans erreur d'outil. L'endpoint officiel r\'epond 404 : \textbf{aucun score n'existe} et aucune valeur z\'ero ne doit \^etre affich\'ee. Seize hachages d'archives sont v\'erifi\'es. L'inspection cibl\'ee des archives pr\'ecise la cause enregistr\'ee : 15 sondes de scan pour 35 n\oe uds d\'eclar\'es (la projection de 36 appareils inclut le runner mais ne d\'emontre pas la couverture des vingt autres n\oe uds), puis deux derni\`eres r\'eponses vides sans outil sur motif \code{length} et arr\^et au tour 5, bien avant tout plafond de 50 tours. Aucune erreur d'outil n'explique cet arr\^et. Le refus de validation pour reconnaissance incompl\`ete est coh\'erent avec ces traces ; le correctif de continuation bornée des réponses sans outil est préparé (section~\ref{sec:corrections}) et reste à valider en laboratoire. Le motif \code{length} ne permet pas, seul, d'attribuer la cause à la taille du contexte. Augmenter un compteur ou accepter une projection incompl\`ete ne r\'esoudrait pas ce d\'efaut.

\section{D\'epassement de p\'erim\`etre et incident d'arr\^et}
\label{sec:scope}
Pendant la campagne, le p\'erim\`etre a \'et\'e r\'eduit \`a S1--S12 avec un observateur demandant l'arr\^et apr\`es S12. L'arr\^et a \'echou\'e : l'observateur POST \code{/api/pipeline/stop} a omis l'authentification administrateur requise et re\c{c}u HTTP 401. Le succ\`es des lectures GET et du streaming ne v\'erifiait pas l'autorisation de mutation. Les démarrages publics de l'API sont exempts d'authentification administrateur, tandis que l'arrêt l'exige. La fin naturelle du lot ne valide pas ce mécanisme d'arrêt.

Cons\'equence : \textbf{tous les sc\'enarios S13--S19 du lot original ont finalement roul\'e} avant le S1 terminal, avec nettoyage, et le lot s'est termin\'e naturellement. Il ne s'agit donc ni d'un arr\^et authentifi\'e r\'eussi ni d'une simple course de quelques secondes du seul S13. Par politique de rapport, ces d\'emarrages et productions hors p\'erim\`etre sont divulgu\'es s\'epar\'ement, leurs archives sont pr\'eserv\'ees, et \textbf{leurs scores \'eventuels sont exclus} de l'analyse S1--S12. Aucun d\'eploiement ne devait avoir lieu pendant le lot actif, et aucun nouveau lot de remplacement n'avait été lancé au gel. Les futurs runs correctifs S1--S12 resteront du d\'eveloppement.

\section{Probl\`emes observ\'es, corrections et validation}
\label{sec:corrections}
\subsection{Regroupement par cause}
Les d\'efauts sont regroup\'es par cause, avec pour chacun les commits li\'es et le statut de validation. Au gel, tous les commits ci-dessous sont \textbf{préparés, poussés, non déployés} ; leur validation r\'eelle sur nato reste attendue. Ils ne modifient ni la v\'erit\'e terrain ni les scores historiques. Le code, les tests logiciels et les rejeux ne prouvent aucun gain.

\textbf{Projection et r\'esum\'es (S2, S3, S6, S1 commun).} Le r\'esum\'e final de phase~4 conservait des compteurs planifi\'es p\'erim\'es apr\`es le suivi de d\'ecouverte, et des red\'ecouvertes d'h\^otes connus faisaient dispara\^itre des v\'erifications planifi\'ees de la projection canonique (S3 \code{VULN-013} ; S6 \code{VULN-014}, \code{VULN-016} ; S1 commun : compteur de candidats 16 contre 17 lignes, ignorés 0 contre 1). Li\'es : \code{eb3b865} (s\'eparation r\'esum\'e/projection ; 177 contr\^oles ; rejeu S2 \`a verdicts et r\'ef\'erences conserv\'es) et \code{f2c4ac6} (exclusion des adresses d\'ej\`a connues ; rejeu S3 conservant 39 tests et 27 confirmations dont \code{VULN-013}, contre 38 et 26 historiques). Validation r\'eelle requise : nouveaux runs S2/S3/S6 et S1 commun apr\`es d\'eploiement.

\textbf{Scanner et sondes redondantes (S2, S3).} N\'egociations WebSocket HTTP tent\'ees apr\`es qu'une sonde nmap fra\^iche et r\'eussie a observ\'e \code{9001/tcp closed} sur un port optionnel non d\'eclar\'e. Li\'e : \code{31501fe} (110 contr\^oles ; cinq rejeux de brokers \`a constats inchang\'es ; les \'ecouteurs ouverts conservent leurs sondes, les ferm\'es omettent les sondes redondantes). Validation r\'eelle : runs S2/S3 apr\`es d\'eploiement, statut scanner et couverture de d\'ecouverte pr\'eserv\'ee.

\textbf{Normalisation des preuves HTTP et robots (S3, S4, S6).} Listes de chemins dans un endpoint singulier rejet\'ees comme non primaires (S3) ; v\'erificateur limit\'e \`a la divulgation de version, aveugle aux chemins publi\'es dans \code{robots.txt} (S4) ; banni\`ere nginx cr\'edit\'ee comme exposition de chemins (S6). Li\'es : \code{44add6c} (premier chemin principal, liste conserv\'ee ; 62 contr\^oles ; rejeu ne changeant que les deux verdicts S3) et \code{ae7abdf} (preuve d\'edi\'ee des chemins \code{robots.txt}, r\'eponse fra\^iche li\'ee \`a la ressource cr\'edit\'ee y compris fichiers secondaires, r\`egles des op\'erations coupl\'ees inchang\'ees ; 474 contr\^oles, 197 interpr\'etations rejou\'ees \`a r\'ef\'erences inchang\'ees). La correction robots conserve le HTTP et exige une preuve propre par ressource. Validation r\'eelle : runs S4/S6 (et S3) apr\`es d\'eploiement ; les r\'esultats devront porter le nouveau contrat et ne seront pas compar\'es \`a contrat constant avec l'historique.

\textbf{MySQL et identifiants (S2, S7, S11).} Mot de passe vide devenu invite interactive, redirection shell pass\'ee comme argument de base, permission syst\`eme inutile exig\'ee ; identifiant \code{root} vide S7 non confirm\'e ; port\'ee MySQL S11 restreinte (pas d'\'ecriture ni d'acc\`es toute adresse). Li\'e (partie MySQL de) \code{eb3b865} (options client documentées~\cite{mysql}, TCP forcé). Validation r\'eelle : runs S2/S7 apr\`es d\'eploiement.

\textbf{D\'ecouverte et port\'ee locale (S4, S2/S7).} Auto-analyse du runner via l'IP cliente MySQL ; inventaire binaire local refus\'e \`a cause de noms de programmes r\'eseau dans la recherche ; composantes du chemin local \code{.ssh/} faussement reconnues comme commandes r\'eseau. La correction de chemin traite \code{.ssh/} dans une arborescence ; la recherche du seul nom \code{.ssh}, sans slash, reste une limite du garde historique. Liés : \code{9d7b285} (exclusion des adresses du runner ; 31 contr\^oles), \code{ba14300} (correction du faux refus ; rejeu de 266 enregistrements ne changeant que ce refus), \code{43e519c} (composantes de chemin ; 50 contr\^oles, rejeu de politique sur 475 enregistrements ne changeant que deux refus, aucune commande ex\'ecut\'ee). Validation r\'eelle : runs S4 (puis S2/S7) apr\`es d\'eploiement ; les refus de port\'ee pour actions r\'eseau restent dus.

\textbf{Alias de service (S4).} Nom Nmap \code{mbap}~\cite{nmap} normalisé mais sans alias Modbus, d'o\`u un diagnostic dupliqu\'e malgr\'e la sonde d\'eclar\'ee. Li\'e : \code{a7136e5} (36 contr\^oles ; fusion sur m\^eme port et transport, diagnostic erron\'e supprim\'e, sans gain de couverture revendiqu\'e). Validation r\'eelle : statut de phase~3 S4 apr\`es d\'eploiement ; le run historique reste inchang\'e.

\textbf{Budgets de temps (S5).} \'Ech\'eances 240~s du routeur et HVAC expos\'ees comme erreur g\'en\'erique avec \code{retry=0}. Li\'e : \code{da6e607} (cause terminale \code{deadline} ; 135 contr\^oles et rejeu des fronti\`eres ; l'observation d'erreur d'origine et la requ\^ete sont pr\'eserv\'ees). Ce correctif classe le diagnostic : il ne rallonge aucun budget et ne garantit pas la conclusion du mod\`ele. Validation r\'eelle : run S5 apr\`es d\'eploiement.

\textbf{Fixtures de laboratoire (S3/S7, S10).} Utilitaire privil\'egi\'e en script Bash SUID, incapable d'impl\'ementer la lecture privil\'egi\'ee d\'eclar\'ee ; fixture Node-RED non pr\^ete (litt\'eral d'octets invalide, contr\^oles de disponibilit\'e non bloquants). Li\'es : \code{d8a9706} (utilitaire compil\'e en lecture seule pour S3/S7 ; 72 contr\^oles, cas Ansible fact/assert uniquement, sans ex\'ecution locale de sc\'enario) et \code{e13d2e1} (correction Node-RED ; 81 contr\^oles logiciels, syntaxes v\'erifi\'ees ; CI \code{37284944781} en succ\`es sur le SHA exact : 3\,159 tests pass\'es, un ignor\'e ; jobs de d\'eploiement ignor\'es). Validation r\'eelle : S3, S7 et S10 sur nato (comportement compil\'e, permissions r\'eelles, disponibilit\'e du service) ; les contr\^oles logiciels ne certifient pas la v\'erit\'e terrain.

\textbf{Continuation de reconnaissance S12.} La reconnaissance compl\`ete retournait pr\'ematur\'ement apr\`es une seconde r\'eponse sans outil sur motif \code{length} alors que la couverture minimale manquait. Li\'e : \code{b39b0b1} (application du comportement strict d\'ej\`a pr\'evu \`a chaque profil de reconnaissance ; 165 tests ciblés ; CI \code{37287309302} en succès ; non déployé au gel). La continuation reste born\'ee : elle n'\'elargit ni le budget mod\`ele/contexte ni les tours, plafonds et arr\^ets inchang\'es, et n'assure pas la convergence du mod\`ele. Validation r\'eelle : run S12 complet en six phases sur nato apr\`es repos/nettoyage et CI du commit exact ; l'\'echec historique reste inchang\'e et aucune observation manquante n'est inf\'er\'ee.

\textbf{Liaison native des ressources (S8--S11).} Endpoint composite FTP bloquant les preuves exactes de lecture de r\'epertoires/fichiers ; lectures natives Redis et SNMP bloqu\'ees par des placeholders HTTP racine ; datagramme SNMP arbitraire accept\'e comme preuve de lecture de communaut\'e et prose candidate p\'erim\'ee ; r\'eponses natives sans bornes de destination, authentification et propri\'et\'e. Li\'e : \code{0f4e5ff} (chemins FTP natifs, lectures Redis typ\'ees, lectures SNMP typ\'ees ; 289 tests ciblés ; rejeu pur de 228 confirmations sur 13 archives hash-v\'erifi\'ees avec exactement 9 modifications d'admission sur S8--S11, entre \code{b39b0b1} préparé/non déployé et \code{0f4e5ff} ; vérité terrain non chargée, scores non recalcul\'es ; \code{git diff --check} pass\'e). Les preuves admises ne sont pas des gains de VP : ne pas inventer de faux n\'egatifs ni de progression. Validation r\'eelle : runs S8--S11 apr\`es d\'eploiement sous contrat \code{strict-v3.16} / \code{evidence-v15}.

\subsection{Statut global de validation des corrections}
Au gel, quatorze commits sont préparés et poussés, tête \code{0f4e5ff} sur \code{main} : les onze premiers avec CI \code{37251387946} en succ\`es sur \code{d8a9706} ; le douzi\`eme Node-RED \code{e13d2e1} avec CI \code{37284944781} en succ\`es ; le treizi\`eme Recon \code{b39b0b1} avec CI \code{37287309302} en succ\`es ; le quatorzi\`eme natif \code{0f4e5ff} avec CI \code{37290434113} \textbf{en cours au gel}, d\'eploiement en attente. Contrainte de d\'eploiement : v\'erifier la CI du commit exact et le SHA r\'eellement d\'eploy\'e avant tout nouveau run ; ne pas d\'eployer pendant un lot actif (condition d\'esormais remplie par la fin naturelle du lot, mais sans d\'eploiement constat\'e au gel). Les validations hors r\'eseau (de 31 \`a 474 contr\^oles cibl\'es selon les commits, plus rejeux d'archives aux empreintes v\'erifi\'ees et 289 tests natifs) documentent explicitement les verdicts conserv\'es et ceux modifi\'es, ainsi que les r\'ef\'erences de preuves. Les scores historiques restent inchang\'es. Ces rejeux ne prouvent ni la finalisation des analyses en laboratoire ni un gain exp\'erimental. \textbf{Aucune validation r\'eelle apr\`es correction n'a eu lieu au gel}, y compris pour le S1 commun qui reste un baseline sous contrat historique.

\paragraph{Suivi postérieur au gel, sans nouveau résultat expérimental.}
La CI finale \code{37290434113} a ensuite réussi sur \code{0f4e5ff882eeb43faed6756a5c0d008a1f6a798d} : 3\,213 tests passés, un ignoré, et déploiement \code{deploy-master / update} réussi. L'API confirme ce SHA déployé et l'état inactif après nettoyage avant le lancement du 5 octobre 2026 à 11:47 (Europe/Brussels). Un lot correctif explicitement borné à S12, S1 puis S2--S11, \code{batch-a4a0a1d88e00}, est lancé avec les six phases, \code{full}, \code{ollama-umons / qwen3.8:27b} et nettoyage automatique. Au moment de cette révision il commence S12 ; aucun résultat de ce lot n'entre dans les tableaux gelés. Les preuves du contrôle et de la requête sont \code{corrective-main-ci.json}, \code{corrective-batch-request.json} et \code{corrective-batch-start-response.json}, sous le dossier d'artefacts. Le contrat attendu de ces futurs résultats est \code{strict-v3.16 / evidence-v15}, à vérifier dans chaque archive.

\section{Discussion et limites}
\label{sec:discussion}
La liaison UMONS fonctionne et des audits complets au sens des six phases produisent des rapports consultables avec nettoyages r\'eussis, mais la campagne reste inachev\'ee sur l'essentiel : r\'esolution de l'\'echec S12, revues s\'emantiques exhaustives S8--S11 et S1 commun d\'etaill\'e, et toutes les validations r\'eelles des quatorze correctifs. La revue conserve le compteur d'acc\`es actuel (preuves d'authentification et de commandes) avec les actions protocolaires Modbus pr\'esent\'ees s\'epar\'ement ; la divulgation de chemins par \code{robots.txt}, mal traitée dans le contrat historique limité aux versions, possède désormais une correction typée préparée dont la validation réelle reste attendue. Les r\'esultats historiques sont conserv\'es pendant cet examen.

Le mod\`ele est stochastique et la reconnaissance varie : les r\'ealisations S1 montrent des F1 de 0,917, 0,846 puis 0,880 sans que l'on puisse conclure \`a une progression. La d\'eduplication conservatrice demeure une limite de qualit\'e des sorties. Les preuves insuffisantes restent ind\'etermin\'ees, les \'echecs de finalisation gardent les observations de scanners, et toute extrapolation au-del\`a des octets et r\'eponses archiv\'es est \'ecart\'ee. Les contr\^oles arithm\'etiques (35 par run avec score) attestent le calcul, pas la s\'emantique ; les exclusions g\'en\'eriques pour preuve insuffisante sur S8--S11 ne sont pas n\'ecessairement des faux n\'egatifs. Les erreurs d'outils \'elev\'ees sur S8--S11 (36, 30, 49, 53) ne sont pas interpr\'et\'ees comme mesure de qualit\'e sans revue des causes.

Limites explicites : aucune agr\'egation des diagnostics avec le lot ; aucun gain scientifique du LLM (ni comparaison avec une automatisation \`a r\`egles ni mesure causale de l'apport du mod\`ele) ; aucune certification de la totalit\'e de la v\'erit\'e terrain (les pr\'e-contr\^oles ne v\'erifient qu'injection de d\'eploiement et playbooks) ; aucune conclusion sur S13--S19, exclus du p\'erim\`etre ; pas de tenu \`a l'\'ecart ind\'ependant. Le rapport est une analyse interm\'ediaire, pas une validation finale de la plateforme. Le contrôle éditorial de cette version ne remplace pas la validation expérimentale ; ses réserves et la portée de l'inspection du PDF sont consignées dans le README.

\section{Conclusion provisoire et suite}
\label{sec:conclusion}
Au gel, le pipeline ex\'ecute les six phases, \'evalue et nettoie sur nato~/~pve-nato, avec des d\'efauts concrets identifi\'es (projection, scanner, normalisation HTTP, MySQL, d\'ecouverte et port\'ee, alias de service, budgets de temps, fixtures, continuation de reconnaissance, liaison native) et quatorze corrections pr\'epar\'ees en attente de validation r\'eelle. Le lot initial est termin\'e : S1 commun 0,880 ; S2--S11 avec scores archiv\'es non r\'e\'evalu\'es (F1 confirm\'es de 0,750 ; 0,545 ; 0,450 ; 0,541 ; 0,686 ; 0,552 ; 0,483 ; 0,538 ; 0,333 ; 0,481) ; S12 en \'echec de reconnaissance sans score ; deux S1 de diagnostic (0,917 ; 0,846) hors lot. Aucune conclusion globale de fiabilit\'e n'est \'etablie.

Prochaines étapes, dans S1--S12 uniquement (développement) : suivre le lot correctif lancé après CI et déploiement vérifiés, puis contrôler les résultats réels de la continuation de reconnaissance S12 ; revalider S2--S11 selon l'effet de chaque correctif (dont S3/S7 pour les fixtures, S10 pour Node-RED, S8--S11 pour le natif), en enregistrant le nouveau contrat \code{strict-v3.16} / \code{evidence-v15} sans comparaison \`a contrat constant et sans r\'e\'ecrire l'\'etape interm\'ediaire \code{3.15/evidence-v14} jamais d\'eploy\'ee ; compl\'eter les revues s\'emantiques exhaustives S8--S11 et S1 commun ; conserver la port\'ee explicite des acc\`es et valider le contrat corrigé de divulgation de chemins sans modifier rétroactivement v\'erit\'e terrain ni scores ; puis figer le bilan avec ses limites.
\clearpage
\phantomsection
\addcontentsline{toc}{section}{R\'ef\'erences}
{\small
\begin{thebibliography}{9}
\raggedright
\setlength{\itemsep}{1pt}
\setlength{\parskip}{0pt}
\bibitem{catalogue} Projet LANCE, \emph{Catalogue et scénarios publics}, \path{benchmarks/catalog.yaml} et \path{docs/benchmark/v1/scenarios.md}, versions inspectées : \code{457d060} et \code{e9d7d55}, octobre 2026.
\bibitem{contrat} Projet LANCE, \emph{Exécution et évaluation du benchmark}, \path{docs/benchmark/v1/execution.md} et \path{docs/benchmark/v1/evaluation.md}, version baseline \code{e9d7d55}, 4 octobre 2026. Historique : \code{strict-v3.14 / evidence-v13}. Code des futurs runs correctifs : \code{0f4e5ff}, \code{strict-v3.16 / evidence-v15} ; étape préparée \code{3.15 / evidence-v14} jamais déployée.
\bibitem{artefacts} Archives opérationnelles, \path{output/validation/2026-10-04-qwen-umons/} (base relative des fichiers ci-après). Faisceau gelé : \path{muse-baseline-report-2026-10-05/authoritative-evidence.json}, snapshot \code{2026-10-05T09:32:42Z}, 14 archives, 1\,582 empreintes. Incident : \path{scope-stop-authentication-incident.json} ; file originale : \path{development-batch-request.json}, \path{development-batch-result.json} ; nouvelle file : \path{corrective-batch-request.json}. Chaque dossier de run possède son \path{collection-manifest.json}.
\bibitem{revues} Sous cette base : \path{s2-provisional-review.json}, \path{s3-provisional-review.json}, \path{s4-provisional-review.json}, \path{s5-phase4-semantic-review.json}, \path{s6-phase4-semantic-review.json}, \path{s7-phase4-semantic-review.json}, \path{s7-fixture-suid-review.json}, \path{s1-common-baseline-review.json}, \path{s8-phase4-semantic-review.json}, \path{s9-phase4-semantic-review.json}, \path{s10-phase4-semantic-review.json}, \path{s11-phase4-semantic-review.json}, \path{native-resource-binding-replay.json}, \path{native-resource-correction.json}, \path{s12-recon-failure-review.json}, \path{s12-recon-continuation-correction.json}. Revue exhaustive des FN et preuves secondaires encore ouverte.
\bibitem{ci} GitHub Actions, validations successives : \url{https://github.com/Tanguyvans/LANCE/actions/runs/37251387946} (\code{d8a9706}, déploiement ignoré), \url{https://github.com/Tanguyvans/LANCE/actions/runs/37284944781} (\code{e13d2e1}, déploiement ignoré), \url{https://github.com/Tanguyvans/LANCE/actions/runs/37287309302} (\code{b39b0b1}, déploiement ignoré), \url{https://github.com/Tanguyvans/LANCE/actions/runs/37290434113} (\code{0f4e5ff}, en cours au gel ; succès et déploiement vérifiés ensuite). Copies : \path{prepared-corrections-ci.json}, \path{nodered-correction-ci.json}, \path{s12-recon-correction-ci.json}, \path{corrective-main-ci.json}, sous la base d'artefacts.
\bibitem{mysql} Oracle, \emph{MySQL 8.0 Reference Manual, mysql Client Options}, options \code{--password} et \code{--skip-password}, \url{https://dev.mysql.com/doc/refman/8.0/en/mysql-command-options.html}, lecture ciblée le 4 octobre 2026.
\bibitem{nmap} Projet Nmap, \emph{nmap-services}, entrée \code{mbap}, TCP 502, \url{https://svn.nmap.org/nmap/nmap-services}, entrée consultée le 4 octobre 2026.
\bibitem{execve} Linux man-pages, \emph{execve(2)}, version 6.19, \url{https://man7.org/linux/man-pages/man2/execve.2.html}, passage sur scripts et bits set-UID/set-GID consulté le 5 octobre 2026.
\bibitem{precedent} Projet LANCE, rapport intermédiaire précédent, révision 5 octobre 2026 10:36, gel 10:12 Europe/Brussels : \path{muse-baseline-report-2026-10-05/previous-report.tex} et \path{previous-report.pdf} dans ce même sous-dossier de la base. Cette version étend le gel à 11:32 sans réécrire les résultats historiques.
\end{thebibliography}
}
\end{document}
